\documentclass{amsart} 

 %margin format 
\usepackage[paperheight=11in, 
    paperwidth=8.5in,
    outer=1.25in,
    inner=1.25in,
    bottom=1.4in,
    top=1.4in]{geometry} 

 \hyphenpenalty = 3000 

 %packages
 %colors package
\usepackage[dvipsnames]{xcolor}
\usepackage{amsmath, amsxtra, amsthm, amsfonts, amssymb}
\usepackage{mathrsfs}
\usepackage{graphicx}
\usepackage{url}
%\usepackage{color}
\usepackage{bbm}
\usepackage{tikz-cd}
\usetikzlibrary{fit}
\renewcommand{\familydefault}{ppl}
\usepackage[T1]{fontenc}
\linespread{1.2}
\usepackage{enumitem}
%floating figures
\usepackage{float}
%Quiver
% `calc` is necessary to draw curved arrows.
\usetikzlibrary{calc}
% `pathmorphing` is necessary to draw squiggly arrows.
\usetikzlibrary{decorations.pathmorphing}
% `spath3` is necessary to draw shortened edges.
\usetikzlibrary{spath3}

% A TikZ style for curved arrows of a fixed height, due to AndréC.
\tikzset{curve/.style={settings={#1},to path={(\tikztostart)
    .. controls ($(\tikztostart)!\pv{pos}!(\tikztotarget)!\pv{height}!270:(\tikztotarget)$)
    and ($(\tikztostart)!1-\pv{pos}!(\tikztotarget)!\pv{height}!270:(\tikztotarget)$)
    .. (\tikztotarget)\tikztonodes}},
    settings/.code={\tikzset{quiver/.cd,#1}
        \def\pv##1{\pgfkeysvalueof{/tikz/quiver/##1}}},
    quiver/.cd,pos/.initial=0.35,height/.initial=0}

% A TikZ style for shortening paths without the poor behaviour of `shorten <' and `shorten >'.
\tikzset{between/.style n args={2}{/tikz/spath/at end path construction={
    \tikzset{spath/split at keep middle={current}{#1}{#2}}
}}}

% TikZ arrowhead/tail styles.
\tikzset{tail reversed/.code={\pgfsetarrowsstart{tikzcd to}}}
\tikzset{2tail/.code={\pgfsetarrowsstart{Implies[reversed]}}}
\tikzset{2tail reversed/.code={\pgfsetarrowsstart{Implies}}}
% TikZ arrow styles.
\tikzset{no body/.style={/tikz/dash pattern=on 0 off 1mm}}

%hyperlink setup 
 \usepackage[hidelinks,backref=page]{hyperref} 
 \hypersetup{
    colorlinks,
    citecolor=magenta,
    filecolor=magenta,
    linkcolor=blue,
    urlcolor=black
}

\usepackage{cleveref}

%letters
\newcommand{\QQ}{\mathbb Q}
\newcommand{\RR}{\mathbb R}
\newcommand{\CC}{\mathbb C}
\newcommand{\PP}{\mathbb P}
\newcommand{\ZZ}{\mathbb Z}
\newcommand{\be}{\mathbf{e}}
\newcommand{\M}{\mathrm{M}}
\newcommand{\kk}{\mathbbm{k}}
\renewcommand{\P}{\mathrm{P}}
\newcommand{\rk}{\operatorname{rk}}

%referenced theorems
\theoremstyle{plain}
\newtheorem{RTh}{Theorem}[section] 
\newtheorem{RLem}[RTh]{Lemma}
\newtheorem{RProp}[RTh]{Proposition}
\newtheorem{RCor}[RTh]{Corollary}

%theorem environments
\theoremstyle{definition}
\newtheorem{claim}{Claim}
\newtheorem{thm}{Theorem}[section]
\newtheorem{cor}[thm]{Corollary}
\newtheorem{lem}[thm]{Lemma}
\newtheorem{prop}[thm]{Proposition}
\newtheorem{defn}[thm]{Definition}
\newtheorem{conj}[thm]{Conjecture}
\newtheorem{eg}[thm]{Example}
\newtheorem{rem}[thm]{Remark}
\newtheorem{obs}[thm]{Observation}
\newtheorem{fact}[thm]{Fact}
\newtheorem{exer}[thm]{Exercise}
\newtheorem{?}[thm]{Question}
\newtheorem{maintheorem}{Theorem}
\renewcommand{\themaintheorem}{\Alph{maintheorem}}

%Symbols
%Commenting Command
\newcommand{\cesay}[1]{\textcolor{cyan}{\textbf{CE says:} #1}}
\newcommand{\dqsay}[1]{\textcolor{BurntOrange}{\textbf{DQ says:} #1}}
\newcommand{\nnsay}[1]{\textcolor{Blue}{\textbf{NN says:} #1}}
%lefts and rights
\def\lang{\left\langle}
\def\rang{\right\rangle}
\def\lnorm{\left\|}
\def\rnorm{\right\|}
\def\labs{\left|}
\def\rabs{\right|}
\def\lbbrace{\{\!\{}
\def\rbbrace{\}\!\}}
\newcommand{\evalat}[2]{\left.#1\right\rvert_{#2}}
\newcommand{\abs}[1]{\left|#1\right|}
%operator names 
\def\conv{\operatorname {conv}}
%stack arrrows
\newcommand{\stackarrow}[2]{\mathrel{\overset{#1}{\underset{#2}{\longrightarrow}}}}


\newcommand{\bx}{\mathbf{x}}
\DeclareMathOperator{\mv}{\mathrm{MV}}


%color
\definecolor{jade}{rgb}{0.0, 0.66, 0.42}
\definecolor{aggieblue}{RGB}{2,40,81}
\definecolor{aggiegold}{RGB}{255,191,0}
\colorlet{aggiegoldT}{aggiegold!50}

\title{Induced Lorentzian polynomials}
\author{Christopher Eur, Nutan Nepal, Daniel Qin}

\begin{document}

\maketitle

\vspace{-10pt}

\begin{abstract}
Given a party of $m$ people, whose expertise collectively covers~$n$ topics, and a subset $T$ of the topics, consider forming a panel of $|T|$ people from the party such that $T$ can be covered by assigning a distinct topic to each panel member with the expertise.
We show that the numbers of such panels, as $T$ varies, form a Lorentzian polynomial.  We achieve this by showing that a certain linear operator on polynomials, which we call the ``inducing operator'' for its connection to induced (poly)matroids, preserves Lorentzian polynomials and realizable volume polynomials.
\end{abstract}


\section{Introduction}

For a positive integer $m$, let $[m] = \{1, \dotsc, m\}$, and let $\mathcal S = (S_1, \dotsc, S_n)$ be a finite sequence of subsets of $[m]$.
For a subset $T \subseteq [n]$, consider the following matching statistic:
\begin{align*}
C_T(\mathcal S) := |\{B \subseteq [m] : \text{there is a bijection $m: T \to B$ with $m(t) \in S_t$ for all $t\in T$}\}|.
\end{align*}
We show that these quantities display the following log-concavity property.  For a nonnegative integer $r$, let $\binom{[n]}{r}$ be the set of cardinality-$r$ subsets of $[n]$, and for $T\subseteq [n]$, write $y^T = \prod_{t\in T} y_t$.

\begin{thm}\label{thm:main1}
For any integer $0\leq r \leq m$, the polynomial $f_{\mathcal S, r} := \sum_{T \in \binom{[n]}{r}}  C_T(\mathcal S) y^T$ is \emph{Lorentzian} in the sense of \cite{BH20}.  More strongly, it is a \emph{realizable volume polynomial} over any field in the sense of \cite{GHMSW2025volpres}.
\end{thm}

See \cite{Huh2026VolumePolynomials} for a survey of Lorentzian polynomials and volume polynomials, as well as inequalities that their coefficients satisfy.
In particular, the theorem implies that the quantities $C_T(\mathcal S)$ satisfy variations of the Alexandrov--Fenchel inequality and the reverse Khovanskii--Teissier inequality.
We do not know elementary proofs of such inequalities for the $C_T(\mathcal S)$.

\smallskip
We caution that a subset $B$ is counted towards $C_T(\mathcal S)$ only once even if the number of bijections $m$ satisfying the condition is greater than one.
As a result, a partial derivative of $f_{\mathcal S,r}$ is usually not another polynomial constructed with a similar matching statistic.
The common inductive strategy involving partial derivatives for showing that a polynomial is Lorentzian therefore does not seem useful in our case.
We deduce the theorem by introducing a linear operator on polynomials that preserve Lorentzian polynomials and realizable volume polynomials, as follows.
Let $\Gamma_{\mathcal S}$ be a bipartite graph on vertices $[m]\sqcup [n]$ where a vertex $i\in [m]$ is connected to a vertex $j\in [n]$ if and only if $i\in S_j$.

\begin{defn}
We say a pair of $\alpha \in \ZZ_{\geq 0}^m$ and $\beta \in \ZZ_{\geq 0}^n$ admits a \emph{$\mathcal S$-matching} if there exists a weighting $w : E(\Gamma_{\mathcal S}) \to \ZZ_{\geq 0}$ of the edges of $\Gamma_{\mathcal S}$ such that, for all $i\in[m] $ and $j\in [n]$, the sum of the weights of the edges incident to $i$ is equal to $\alpha_i$, and to $j$ is equal to $\beta_j$.
We write $\alpha\overset{\mathcal S}{\longleftrightarrow}\beta$ in this case.
\end{defn}

For $\alpha \in \ZZ_{\geq 0}^m$, write $\alpha! = \alpha_1 ! \dotsb \alpha_m!$ and $x^\alpha = x_1^{\alpha_1}\dotsb x_m^{\alpha_m}$.  

\begin{thm}\label{thm:main2}
The linear operator $\mathcal I_{\mathcal S}: \RR[x_1, \dotsc, x_m] \to \RR[y_1, \dotsc, y_n]$ defined on the monomials by 
\[
\mathcal I_{\mathcal S}\left(\frac{x^\alpha}{\alpha!}\right) := \sum_{\beta: \alpha\overset{\mathcal S}\longleftrightarrow\beta} \frac{y^\beta}{\beta!}
\]
preserves Lorentzian polynomials and realizable volume polynomials over a field.
\end{thm}

The operator $\mathcal I_{\mathcal S}$ does not always preserve \emph{stable polynomials} in the sense of \cite{BB09}; see Example~\ref{eg:notstable}.
We prove Theorem~\ref{thm:main2} in Section~\ref{sec:symbol} by showing that the symbol of the operator $\mathcal I_{\mathcal S}$ is a realizable volume polynomial over any field.
%, and deduce Theorem~\ref{thm:main1} as a corollary.

\begin{proof}[Proof of Theorem~\ref{thm:main1} from Theorem~\ref{thm:main2}]
The polynomial $\sum_{T \in \binom{[n]}{r}}  C_T(\mathcal S) y^T$ is the result of taking the multi-affine part (i.e.\ the terms with square-free monomials) of the polynomial $\mathcal I_{\mathcal S}(\sum_{B\in \binom{[m]}{r}} x_B)$.
The elementary symmetric polynomial $\sum_{B\in \binom{[m]}{r}} x_B$ is a realizable volume polynomial over any field because it is over any infinite field $\kk$, witnessed by the arrangement Schubert variety of a generic $r$-dimensional subspace of $\kk^m$ \cite[Example 3.5]{Huh2026VolumePolynomials}, and taking field extensions do not change the set of realizable volume polynomials \cite[Proposition 2.10]{GHMSW2025volpres}.
The desired conclusion now follows from Theorem~\ref{thm:main2} and that taking the multi-affine part preserves Lorentzian polynomials and realizable volume polynomials \cite[Corollary~3.5]{BH20}, \cite[Corollary~3.3]{GHMSW2025volpres}.
\end{proof}

We call the operator $\mathcal I_{\mathcal S}$ the \emph{inducing operator} of $\mathcal S$ because we view it as the polynomial analogue of the construction in matroid theory called \emph{induced (poly)matroids}, as we now explain.
A \emph{polymatroid} $\P$ on $[m]$ is a {function} $\operatorname{rk}_{\P}: 2^{[m]} \to \ZZ_{\geq 0}$ satisfying
\begin{enumerate}[label = (\roman*)]
\item $\operatorname{rk}_\P(\emptyset) = 0$,
\item if $I\subseteq J \subseteq [m]$ then $\operatorname{rk}_\P(I) \leq \operatorname{rk}_\P(J)$, and
\item for all subsets $I$ and $J$ of $[m]$, one has $\operatorname{rk}_\P(I) + \operatorname{rk}_\P(J) \geq \operatorname{rk}_\P(I\cup J) + \operatorname{rk}_\P(I\cap J)$.
\end{enumerate}
We call $\operatorname{rk}_\P$ the \emph{rank function} of $\P$, and define the \emph{base polytope} $Q(\P)$ of $\P$ by
\[
Q(\P) := \Big\{x \in \RR^m : \sum_{i\in I} x_i \leq \operatorname{rk}_\P(I) \text{ for all $I\subseteq [m]$}\quad \text{and}\quad \sum_{i\in [m]} x_i = \operatorname{rk}_\P([m])\Big\}.
\]
A \emph{matroid} on $[m]$ is a polymatroid whose rank function further satisfies $\operatorname{rk}_\P(I) \leq |I|$ for all $I \subseteq [m]$.
We point to \cite{Welsh} as a standard reference on matroids and polymatroids.
Following \cite[Chapter~8.2]{Welsh}, from a polymatroid $\P$ on $[m]$ and a sequence of subsets $\mathcal S = (S_1, \dotsc, S_n)$ of $[m]$, one can construct a new polymatroid $\mathcal I_{\mathcal S}(\P)$ on $[n]$, called the \emph{induced polymatroid}, by setting its rank function to be $\operatorname{rk}_{\mathcal I_{\mathcal S}(\P)}(I) := \operatorname{rk}_\P(\bigcup_{i\in I}S_i)$ for all $I\subseteq [n]$.  By setting the rank function to be $I\mapsto \min\{|I|, \operatorname{rk}_\P(\bigcup_{i\in I}S_i)\}$ instead, one obtains the \emph{induced matroid}.

\smallskip
The operator $\mathcal I_{\mathcal S}$ gives rise to induced polymatroids as follows.
If $f  = \sum_{\alpha} c_\alpha x^\alpha \in \RR[x_1, \dotsc, x_m]$ is Lorentzian, then its \emph{support}, i.e.\ $\operatorname{supp}(f) = \{\alpha : c_\alpha \neq 0\}$, is always the set $Q(\P) \cap \ZZ^m$ of the lattice points of the base polytope of some polymatroid $\P$ \cite[Theorem~2.25]{BH20}.
We say simply that the support of $f$ is the polymatroid $\P$ in this case.
The Hall--Rado theorem for polymatroids, recalled here as Proposition~\ref{prop:HR}, implies that 
\begin{equation*}\tag{$*$}\label{eq:induced}
\text{$\operatorname{supp}(f) = Q(\P) \cap \ZZ^m$ $\implies$ $\operatorname{supp}(\mathcal I_{\mathcal S}(f)) = Q(\mathcal I_{\mathcal S}(\P)) \cap \ZZ^n$}.
\end{equation*}
The operator $\mathcal I_{\mathcal S}$ is not the only operator satisfying the property \eqref{eq:induced} however; for example, the change of variables $x_i \mapsto \sum_{j: S_j \ni i} y_j$ also has the property \eqref{eq:induced}.
We point out how $\mathcal I_{\mathcal S}$ is the "center point" among these in Proposition~\ref{prop:1param}.

\subsection*{Acknowledgements}
We thank American Institute of Mathematics for hosting the workshop ``The geometry of polynomials in combinatorics and sampling'' during which the first author proposed this project.  The first author is supported by US NSF grant DMS-2246518.\textcolor{red}{Add more.}


\section{The symbol of $\mathcal I_{\mathcal S}$}\label{sec:symbol}

We use the ``symbol theory'' originally developed in \cite{BB09} in the context of stable polynomials to prove Theorem~\ref{thm:main2}.
For $\kappa \in \ZZ_{\geq 0}^m$, let $\RR[x_1, \dotsc, x_m]_\kappa$ denote the set of polynomials whose support is contained in the box $[0,\kappa_1] \times \dotsb \times [0, \kappa_m]$.  For $\alpha\in \ZZ_{\geq 0}^m$ we write $\alpha \leq \kappa$ if $\alpha_i \leq \kappa_i$ for all $i \in [m]$.  The \emph{symbol} of a linear operator $T: \RR[x_1, \dotsc, x_m]_\kappa \to \RR[y_1, \dotsc, y_n]$ is defined as
\[
sym_T(y,u) := T\big((x_1+u_1)^{\kappa_1}\dotsb (x_m+u_m)^{\kappa_m}\big) \in \RR[y_1,\ldots,y_n,u_1,\ldots,u_m]
\]
where $T$ acts trivially on $u$, or equivalently, $sym_T(y,u) = \sum_{0\leq \alpha\leq  \kappa} \binom{\kappa}{\alpha}T(x^\alpha) u^{\kappa - \alpha}$ where $\binom{\kappa}{\alpha} := \binom{\kappa_1}{\alpha_1} \dotsb \binom{\kappa_m}{\alpha_m}$.

\begin{prop}\label{prop:symbol}
Let $T$ be a linear operator $\RR[x_1, \dotsc, x_m]_\kappa\to \RR[y_1, \dotsc, y_n]$ that preserves homogeneous polynomials.
\begin{itemize}
\item \cite[Theorem~3.2]{BH20} If the symbol of $T$ is Lorentzian, then $T$ preserves Lorentzian polynomials.
\item \cite[Theorem~1.12]{GHMSW2025volpres} If the symbol of $T$ is a realizable volume polynomial over a field $\kk$, then $T$ preserves realizable volume polynomials over $\kk$.
\end{itemize}
\end{prop}

We now fix an arbitrary $\kappa \in \ZZ_{\geq 0}^m$ for the rest of the section, and consider $\mathcal I_{\mathcal S}$ as an operator on $\RR[x_1, \dotsc, x_m]_\kappa$.
We compute the symbol of $\mathcal I_{\mathcal S}$ to be
\[
sym_{\mathcal I_{\mathcal S}}(y,u) = \sum_{0\leq \alpha\leq  \kappa} \binom{\kappa}{\alpha}\mathcal I_{\mathcal S}(x^\alpha) u^{\kappa - \alpha}
= \kappa! \sum_{0\leq \alpha\leq \kappa} \mathcal I_{\mathcal S}\left(\frac{x^\alpha}{\alpha!}\right) \frac{u^{\kappa-\alpha}}{(\kappa-\alpha)!}
= \sum_{0\leq \alpha \leq \kappa} \sum_{\beta: \alpha\overset{\mathcal S}\longleftrightarrow\beta} \frac{y^\beta}{\beta!}\frac{u^{\kappa-\alpha}}{(\kappa-\alpha)!}.
\]
Note that the symbol is the exponential generating function of its support, which is necessarily a polymatroid if the symbol is to be Lorentzian.
Our primary goal concerns the linear realizability of this polymatroid.
A polymatroid $\P$ on $[m]$ is said to be \emph{linearly realizable} over a field $\kk$ if there exist $\kk$-vector spaces $V_1, \dotsc, V_m$ and a subspace $L \subseteq \bigoplus_{i\in [m]} V_i$ such that
\[
\operatorname{rk}_\P(A) = \dim \Big(\text{image of $L$ under the projection $\bigoplus_{i\in [m]} V_i \to \bigoplus_{j\in A} V_j$}\Big) \quad\text{for all $A\subseteq [m]$}.
\]


\begin{thm}\label{thm:symbol}
The symbol $sym_{\mathcal I_{\mathcal S}}$ of $\mathcal I_{\mathcal S}$ is the exponential generating function of the lattice points of the base polytope a polymatroid on $[m]\sqcup [n]$ that is linearly realizable over any infinite field.
%In particular, $sym_{\mathcal I_{\mathcal S}}$ is Lorentzian, and is a realizable volume polynomial over any field.
\end{thm}

This theorem implies our main result Theorem~\ref{thm:main2} as follows.

\begin{proof}[Proof of Theorem~\ref{thm:main2}]
\cite[Corollary~1.4 \& Remark~3.21]{EL2024intersection} states that the exponential generating function of a linearly realizable polymatroid over $\kk$ is a realizable volume polynomial over $\kk$.
Theorem~\ref{thm:symbol} thus implies that $sym_{\mathcal I_S}$ is a realizable volume polynomial over any infinite field, and in fact over any field because the set of realizable volume polynomials do not change under field extensions \cite[Proposition 2.10]{GHMSW2025volpres}.  Realizable volume polynomials are also Lorentzian \cite[Theorem~4.6]{BH20}.  Proposition~\ref{prop:symbol} now implies the theorem.
\end{proof}

We prepare for the proof of Theorem~\ref{thm:symbol} with the following variant of the Hall--Rado theorem for polymatroids.

\begin{prop}\label{prop:HR}\cite[Theorem 2]{McD75}
Let $\P$ be a polymatroid on a finite set $E$, and let $\mathcal E = (E_1, \dotsc, E_\ell)$ be a sequence of subsets of $E$.
Then, for $\delta \in \ZZ_{\geq 0}^\ell$, we have that $\delta \in Q(\mathcal I_{\mathcal E}(\P))$ if and only if there exists $\gamma\in \ZZ_{\geq 0}^E$ such that the pair $(\gamma,\delta) \in \ZZ_{\geq 0}^E \times \ZZ_{\geq 0}^\ell$ admits a $\mathcal E$-matching.
\end{prop}

\begin{proof}[Proof of Theorem~\ref{thm:symbol}]
We first identify the polymatroid.
For an integer $r$ and a finite set $[N]$, let $\Delta_{[N]}^r$ be the polymatroid on $[N]$ whose rank function is $\rk_\P(A) = r$ for all nonempty $A\subseteq [N]$, or equivalently, its base polytope is $Q(\Delta^r_{[N]}) = \{ x\in \RR_{\geq 0}^{[N]} : \sum_{i\in {[N]}} x_i = r\}$.  It is linearly realizable over any infinite field $\kk$ by a generic embedding of $\kk^r$ in $(\kk^r)^{\oplus [N]}$.
Let $E$ be the set of edges of the bipartite graph $\Gamma_{\mathcal S}$ on $[m]\sqcup [n]$, and for each vertex $v\in [m]\sqcup [n]$, let $E_v$ be the set of edges incident to $v$.
Let $\P$ be the polymatroid on $E$ defined as the direct sum $\Delta^{\kappa_1}_{E_{i\in [m]}} \oplus \dotsb \oplus \Delta^{\kappa_m}_{E_{m\in [m]}}$.
Here we wrote $E_{i\in[m]}$ for $E_i$ to avoid possible confusion with $E_{i\in [n]}$.
Lastly, let $\mathcal E$ be the sequence $(E_v)_{v\in [m]\sqcup [n]}$, and consider the induced polymatroid $\mathcal I_{\mathcal E}(\P)$ on $[m]\sqcup [n]$, which is linearly realizable over any infinite field $\kk$ because $\P$ is.


We claim that a pair $(\alpha,\beta)\in \ZZ_{\geq 0}^m \times \ZZ_{\geq 0}^n$ with $0\leq \alpha\leq \kappa$ admits a $\mathcal S$-matching if and only if $(\kappa-\alpha, \beta) \in Q(\mathcal I_{\mathcal E}(\P))$.
Since $sym_{\mathcal I_{\mathcal S}}$ is the exponential generating function of the set of such $(\kappa-\alpha,\beta)$, the theorem follows.
Proposition~\ref{prop:HR} implies that $(\kappa-\alpha, \beta) \in Q(\mathcal I_{\mathcal E}(\P))$ if and only if there exists $\gamma \in Q(\P) \subset \ZZ_{\geq 0}^E$ such that $\gamma$ admits a $\mathcal E$-matching to $(\kappa-\alpha, \beta)$.
%Since every $e\in E$, considered as a vertex in the bipartite graph $\Gamma_{\mathcal E}$, is adjacent to two other vertices, one from each $[m]$ and $[n]$,
The edge weighting of $\Gamma_{\mathcal E}$ witnessing such a $\mathcal E$-matching can be viewed as a function $\mathbf w: E \to \ZZ_{\geq 0}^2$, $e \mapsto (w_e, w'_e)$ on the edges of $\Gamma_{\mathcal S}$ such that:
\begin{itemize}
\item $(w_e + w'_e)_{e\in E} \in Q(\P)$, 
\item for every $i\in [m]$, the sum of $w_e$ over all $e\in E$ incident to $i$ is $\kappa_i-\alpha_i$, and
%$\displaystyle \sum_{e\in E_{i\in [m]}} w_e = \kappa_i - \alpha_i$ for all $i\in [m]$, and
\item for every $j\in [n]$, the sum of $w'_e$ over all $e\in E$ incident to $j$ is $\beta_j$.
%$\displaystyle \sum_{e\in E_{j\in [n]}} w_e = \beta_j$ for all $j\in [n]$.
\end{itemize}
Since $\P$ is the direct sum the $\Delta^{\kappa_i}_{E_{i\in [m]}}$, we have $Q(\P) = \prod_{i\in [m]} Q(\Delta^{\kappa_i}_{E_{i\in [m]}})$, and thus $(w_e + w'_e)_{e\in E} \in Q(\P)$ if and only if $\sum_{e\in E_{i\in [m]}} (w_e + w'_e) = \kappa_i$ for every $i\in [m]$.  Hence, the last two properties of $\mathbf w$ imply that the edge weighting $w': E \to \ZZ_{\geq 0}$, $e\mapsto w'_e$ witnesses that $(\alpha,\beta)$ admits a $\mathcal S$-matching.
Conversely, if $(\alpha,\beta)$ admits a $\mathcal S$-matching witnessed by an edge weighting $w': E \to \ZZ_{\geq 0}$, $e\mapsto w'_e$, then we may extend $w'$ to a function $\mathbf w: E \to \ZZ_{\geq 0}^2$ satisfying the three bulleted properties above by setting $(w_e)_{e\in E}$ to be any nonnegative integer vector satisfying the second bulleted property.
\end{proof}

\begin{rem}
Let $E$ be the edges of the bipartite graph $\Gamma_{\mathcal S}$.
Given a function $r: E \to \ZZ_{\geq 0}$, we say that a pair $(\alpha, \beta)\in \ZZ_{\geq 0}^m \times \ZZ_{\geq 0}^n$ admits a \emph{$r$-restricted $\mathcal S$-matching} if it admits a $\mathcal S$-matching witnessed by an edge weighting $w: E\to \ZZ_{\geq 0}$ satisfying $w_e \leq r_e$ for all $e\in E$.  For any $r$ and $\kappa$, the proof of Theorem~\ref{thm:symbol} modifies straightforwardly to show that the symbol of the analogue of the operator $\mathcal I_{\mathcal S}$ for $r$-restricted $\mathcal S$-matchings is also a realizable volume polynomial over any field.  The only required change is to replace the polymatroid $\Delta^{\kappa_i}_{E_{i\in [m]}}$ by the polymatroid that is linearly realizable by a generic embedding of $\kk^{\kappa_i}$ in $\bigoplus_{e\in E_{i\in [m]}} \kk^{r_e}$.
\end{rem}

A polynomial $f(x_1, \dotsc, x_n)$ is \emph{stable} if it has no roots in the upper half-plane of $\CC^n$.  Exponential generating functions of polymatroids are often not stable.  In agreement, the operator $\mathcal I_{\mathcal S}$ often does not preserve stability.

\begin{eg}\label{eg:notstable}
Let $\mathcal{S}=(123,123,12)$ be a sequence of subsets of $[3]$.
By noting that every $\beta \in \ZZ_{\geq 0}^3$ with $\beta_1 + \beta_2 + \beta_3 = 3$  except $(0,0,3)$ admits a $\mathcal S$-matching to $(1,1,1)$, one computes that
\[
\mathcal I_{\mathcal S}(x_1x_2x_3) = \frac{1}{6}((y_1+y_2+y_3)^3-y_3^3).\]
We have a factorization
$$(y_1+y_2+y_3)^3-y_3^3=(y_1+y_2)((y_1+y_2+y_3)^2+(y_1+y_2+y_3)y_3+y_3^2)=(y_1+y_2)y_3^2(v^2+v+1)$$ where $v = (y_1+y_2+y_3)/y_3.$
Note that the roots of $v^2+v+1$ lie in the upper half plane of $\CC$.  Using this, we can construct a solution to $I_{\mathcal S}(x_1x_2x_3) = 0$ such as $y_3 = e^\frac{\pi i}{12}$ and $y_1=y_2=\frac{1}{2}(e^{\frac{3\pi i}{4}} - e^{\frac{\pi i}{12}})$, which witnesses that $\mathcal I_{\mathcal S}(x_1x_2x_3)$ is not stable.
\end{eg}



\begin{rem}
We note the following generalization of Theorem~\ref{thm:main1}, which can be deduced from Theorem~\ref{thm:main2} similarly.  For a matroid $\M$ on $[m]$ of rank $r$, and a subset $T\subseteq [n]$, let
\begin{align*}
C_T(\M,\mathcal S) := |\{B \text{ a basis of }\M : \text{there is a bijection $m: T \to B$ with $m(t) \in S_t$ for all $t\in T$}\}|.
\end{align*}
Then the polynomial $\sum_{T\in \binom{[n]}{r}} C_T(\M,\mathcal S) y^T$ is Lorentzian, and further a realizable polynomial over a field $\kk$ if the matroid $\M$ is realizable over a field of the same characteristic.
\end{rem}




\section{Inducing operators}

Proposition~\ref{prop:HR} implies that the linear operator has the following support property \eqref{eq:induced}:
\begin{quote}
If the support of a Lorentzian polynomial $f\in \RR[x_1, \dotsc, x_m]$ is a polymatroid $\P$, then the support of the polynomial $\mathcal I_{\mathcal S}(f)$ is the induced polymatroid $\mathcal I_{\mathcal S}(\P)$.
\end{quote}
There are other linear operators with this property \eqref{eq:induced}, such as the following.


\smallskip
Let $A_{\mathcal S}$ be the matrix whose columns are $\be_{S_1},\dotsc, \be_{S_n}$, where $\be_S := \sum_{i\in S} \be_i \in \RR^m$.  The linear operator $A_{\mathcal S}: f(x) \mapsto f(A_{\mathcal S}y)$, or equivalently the change of variables operator $x_i \mapsto \sum_{S_j \ni i} y_j$, satisfies the support property \eqref{eq:induced}.  Indeed, it is straightforward to verify that the support of $A_{\mathcal S}(x^\alpha)$ and that of $\mathcal I_{\mathcal S}(x^\alpha)$ coincide for $\alpha \in \ZZ_{\geq 0}^m$.
Combinatorially, the distinction between $\mathcal I_{\mathcal S}$ and $A_{\mathcal S}$ is that for $A_{\mathcal S}$, the coefficient of $y^\beta/\beta!$ in $A_{\mathcal S}(x^\alpha/\alpha!)$ is given by the number of edge weightings of $\Gamma_{\mathcal S}$ witnessing that $\alpha \overset{\mathcal S}\longleftrightarrow \beta$, where $\mathcal I_{\mathcal S}$ only detects the existence of such an edge weighting.

\smallskip
More generally, let $A$ be a nonnegative $m\times n$ matrix whose zero-nonzero pattern is the same as that of the matrix whose columns are $\be_{S_1},\dotsc, \be_{S_n}$, where $\be_S := \sum_{i\in S} \be_i \in \RR^m$.  Then the linear operator $A: f(x) \mapsto f(Ay)$ satisfies the support property \eqref{eq:induced}.
The symbol of $A$ is Lorentzian, being a product of nonnegative linear forms.

\begin{prop}\label{prop:1param}
Fix any $\kappa\in \ZZ_{\geq 0}^m$.  Let $\mathbf I_{\mathcal S}$ be the space of linear operators $T: \RR[x]_\kappa \to \RR[y]$ with Lorentzian symbol and with property \eqref{eq:induced}.  Then $\mathbf I_{\mathcal S}$ deformation retracts to the point $\mathcal I_{\mathcal S}$.
\end{prop}

\begin{proof}
The symbol of such an operator $T\in \mathbf I_{\mathcal S}$ is of the form
\[
\sum_{\substack{\alpha,\beta: 0\leq \alpha\leq \kappa\\ \text{and } \alpha \overset{\mathcal S}\longleftrightarrow \beta}} c_{\alpha,\beta} \frac{y^\beta}{\beta!} \frac{u^{\kappa-\alpha}}{(\kappa-\alpha)!}
\]
for some $c_{\alpha,\beta} > 0$.  For any $0\leq q\leq 1$, define $T^q$ to be the operator whose symbol is
\[
\sum_{\substack{\alpha,\beta: 0\leq \alpha\leq \kappa\\ \text{and } \alpha \overset{\mathcal S}\longleftrightarrow \beta}} c_{\alpha,\beta}^q \frac{y^\beta}{\beta!} \frac{u^{\kappa-\alpha}}{(\kappa-\alpha)!}.
\]
We have $T^0 = \mathcal I_{\mathcal S}$.  For $0\leq q\leq 1$, taking $q$-th power of normalized coefficients of a Lorentzian polynomial yields a Lorentzian polynomial \cite[Proposition~3.25]{BH20}, so $T^q \in \mathbf I_{\mathcal S}$ for all $0\leq q\leq 1$.
\end{proof}

Since the symbol of $A_{\mathcal S}$ is Lorentzian, the proof of the proposition provides a quick way to verify that the symbol of $\mathcal I_{\mathcal S}$ is Lorentzian.  However, while the symbol of $A_{\mathcal S}$ is a realizable volume polynomial (over any field), being a product of nonnegative linear forms, one cannot conclude similarly that $sym_{\mathcal I_{\mathcal S}}$ is a realizable volume polynomial.  In fact, we ask:

\begin{?}
For a polynomial $f = \sum_\alpha c_\alpha \frac{x^\alpha}{\alpha}$, denote $f^q := \sum_\alpha c_\alpha^q \frac{x^\alpha}{\alpha}$.  If $f^1$ is a (realizable) volume polynomial (over a field $\kk$), is $f^q$ a volume polynomial (over $\kk$) for all $0\leq q \leq 1$?
\end{?} 
 
One may seek a different 1-parameter family of operators in $\mathbf I_{\mathcal S}$, consisting of operators whose symbols are volume polynomials, that connects $\mathcal I_{\mathcal S}$ and $A_{\mathcal S}$.  The following example shows that this cannot be achieved by varying the coefficients of the matrix $A$.

\begin{eg}
Let $\mathcal S = (1,2,12)$ be a sequence of subsets of $[2]$.  We have $\mathcal I_{\mathcal S}(x_1x_2) = y_1y_2 + y_1y_3 + y_2y_3 + \frac{y_3^2}{2}$, whereas $A(x_1x_2)$ is of the form $ (ay_1 + by_3)(cy_2+dy_3) = acy_1y_2 + ady_1y_3 + bcy_2y_3 + bdy_3^2$.  For all values of $(a,b,c,d)$, the latter has the property that the product of the coefficients of $y_1y_2$ and $y_3^2$ is equal to the product of the coefficients of $y_1y_3$ and $y_2y_3$.
\end{eg}

The desired family of operators can be achieved using Theorem~\ref{thm:symbol} as follows.
For $\mathcal{S}=(S_1,\ldots,S_n)$ on $[m]$, say $S_i=\left\{j_1^{(i)},\ldots,j_{N_i}^{(i)}\right\}$, define another sequence $\widetilde{\mathcal S}$ of subsets of $[m]$ by
$$\widetilde{\mathcal{S}}=(S_1,\ldots,S_n;j_1^{(1)},\ldots,j_{N_1}^{(1)},\ldots,j_1^{(n)},\ldots,j_{N_n}^{(n)}).$$
Let $N = \sum_{i = 1}^nN_i$.  By Theorem~\ref{thm:symbol}, the symbol $sym_{\mathcal I_{\widetilde{\mathcal S}}}(y_1, \dotsc, y_n, \widetilde y_1, \dotsc, \widetilde y_N, u_1, \dotsc, u_m)$ is a realizable volume polynomial over any field $\kk$.
Let $\boldsymbol y'$ be the vector $(y_1, \dotsc, y_1, y_2, \dotsc, y_2, \dotsc, y_n, \dotsc, y_n)$ where $y_i$ is repeated $N_i$ many times.
For any nonnegative vectors $\boldsymbol a\in \RR^n_{\geq 0}$ and $\boldsymbol b\in \RR_{\geq 0}^N$, the nonnegative linear change of variables $sym_{\mathcal I_{\widetilde{\mathcal S}}}(\boldsymbol a \cdot \boldsymbol y, \boldsymbol b \cdot \boldsymbol{ y'}, \boldsymbol u)$ is a volume polynomial over any field.  Let $T_{\boldsymbol a, \boldsymbol b}$ be the linear operator with symbol $sym_{\mathcal I_{\widetilde{\mathcal S}}}(\boldsymbol a \cdot \boldsymbol y, \boldsymbol b \cdot \boldsymbol{ y'}, \boldsymbol u)$. We find that $T_{\mathbf 1,0} = \mathcal I_{\mathcal S}$ and $T_{0,\mathbf 1} = A_{\mathcal S}$.  Thus, any path $(\mathbf{1},0)\rightsquigarrow(0,\mathbf{1})\subseteq \RR_{\geq 0}^n\times \RR_{\geq 0}^{N}$ yields the desired 1-parameter family.


\begin{rem}
In \cite{KMT}, the authors study the analogue of the induced polymatroids construction for M-convex functions.  There, the key strategy is that every bipartite graph can be considered as a composition of simple bipartite graphs, called elementary splits and elementary aggregations.  The same strategy cannot be adopted establish the properties of $\mathcal I_{\mathcal S}$ because the composition of the operators of the form $\mathcal I_{\mathcal S}$ does not respect composition of bipartite graphs.
\end{rem}

{\fontsize{10.5}{11}\selectfont \bibliographystyle{plain}
  \bibliography{references}  }
\end{document}